
\section{Methods}
\label{s:methods}

\subsection{Governing equations}

We perform numerical simulations of convection in deformed spherical shell geometries rotating with constant rotation vector $\boldsymbol{\Omega} = \Omega \, \hz$. The non-dimensional inner and outer radii are denoted $r_i$ and $r_o$, respectively. In all cases presented the aspect ratio of the geometry is fixed to the core-like value of $\chi = r_i/r_o = 0.35$, such that $r_i = \chi/(1-\chi) = 7/13$ and $r_o = 1/(1-\chi) = 20/13$. The inner boundary is spherical in all cases, whereas the outer boundary is deformed according to the relationship
\be
r_{cmb} = r_o - \epsilon  h\lb \theta,\phi\rb,
\label{e:rcmb}
\ee
where $h\lb \theta,\phi\rb$ is the topographic shape function and $\epsilon$ is the maximum topographic amplitude. The colatitude and azimuthal angle are $\theta$ and $\phi,$ respectively. In a co-rotating frame the non-dimensional governing equations for a Boussinesq fluid are
given by
\begin{subequations}
\begin{equation}
	\frac{\partial \mathbf{u}}{\partial t} + \mathbf{u}\cdot \nabla \mathbf{u} + \frac{2}{Ek}\hat{\mathbf{z}} \times \mathbf{u} = -\nabla p + \frac{\Ra}{Pr}\frac{\mathbf{r}}{r_{o}} \vartheta + \nabla^{2}\mathbf{u} ,
	\label{e:mom}
\end{equation}
\begin{equation}
	\frac{\partial \vartheta}{\partial t} + \mathbf{u}\cdot \nabla \vartheta = \frac{1}{Pr }\nabla^{2}\vartheta ,
	\label{e:energy}
\end{equation}
\begin{equation}
	\nabla\cdot\mathbf{u} = 0 ,
	\label{e:incom}
\end{equation}
\label{e:gov_eqn}
\end{subequations}
where the velocity, pressure and temperature are denoted by $\ub $, $p$, and $\vartheta$, respectively, and the position vector is $\mathbf{r}$. The equations have been non-dimensionalised using the shell depth $d$, viscous diffusion time $d^2/\nu$, and temperature $\Delta T$, which is the temperature difference between the inner and outer boundaries.
We use $u,v,$ and $w$ to refer to the $\hat{\mathbf{x}},\hat{\mathbf{y}}$, and $\hat{\mathbf{z}}$ components of the velocity field respectively, and $u_{r}, u_{s},$ and $u_{\phi}$ to refer to the
$\hat{\mathbf{r}},\hat{\mathbf{s}}$, and $\hat{\boldsymbol{\phi}}$ components, where $s$ is the radial coordinate in cylindrical coordinates. 
%The second term on the right hand side of equation (\ref{e:gov_eqn}a) is the buoyancy term, which we abbreviate as $\mathbf{F_b}.$

The non-dimensional control parameters appearing above are the Ekman, Rayleigh, and Prandtl numbers defined, respectively, as
\be
Ek = \frac{\nu}{\Omega d^{2}},\;\; Ra = \frac{g_o \alpha \Delta T d^{3}}{\nu\kappa}, \;\; Pr = \frac{\nu}{\kappa},
\ee
where $\nu$ is the kinematic viscosity, $g_o$ is the gravitational acceleration at the CMB, $\alpha$ is the coefficient of thermal expansion, and $\kappa$ is the thermal diffusivity of the fluid. 


The primary goal of the present study is to determine the role played by the topographic parameters $\epsilon $ and $h\lb \theta,\phi\rb $. Thus, to narrow the parameter space, the Prandtl number is fixed at $Pr = 1$. 
With the exception of five simulations performed at $Ek = 10^{-4},$ the Ekman number is fixed at $Ek = 10^{-6}.$ 
No-slip, fixed temperature boundary conditions are employed. The temperature at $r_{cmb}$ is set to zero in all simulations except for a small set of laminar simulations which we discuss in if the first section of Results. 
In the absence of topography, the onset of convection occurs at the critical Rayleigh number $Ra_c \approx 1.8\times 10^8$ and critical azimuthal wavenumber $m_c = 32$ \citep[e.g.][]{aB22}. 
Because the critical Rayleigh number increases with decreasing $Ek,$ it is useful to define the reduced Rayleigh number
$\Rat = Ra Ek^{4/3},$ 
which is an $O\lb 1\rb $ quantity as $Ek\rightarrow0$ \citep[e.g.][]{eD04}. For convection in a perfect spherical shell of aspect ratio $\chi = 0.35$, the reduced critical Rayleigh number is $\Rat_{c} \approx 1.8$. 
However, as mentioned previously, unique instabilities arise in the deformed geometries which yield non-zero flows for Rayleigh numbers below this value. 

The results in this study are grouped into two categories. We first present results for small values of $\Rat,$ such that $\Rat<\Rat_{c}.$
As we show, some of these cases are unstable to fluid motion due to topographic instabilities. We will nevertheless refer to these simulations as ``subcritical'' because $\Rat<Ra_{c}.$
The second group of results are turbulent, and for these cases we fix $\Rat = 40$ and $Ek = 10^{-6}$. Tables \ref{t:p_space} - \ref{t:p_space3} in the Appendix provide a comprehensive list of the run parameters.


\subsection{Numerical Methods}

The simulations are performed with the spectral element code \texttt{Nek5000} \citep{nek5000,pA84}. We use a cubed sphere to generate the mesh. In the radial direction we use $24$ elements distributed on a Chebyshev grid. This corresponds to a radial resolution of $192$ points.
The longitudinal and latitudinal dimensions are composed of $6$ panels which are $40\times 40$ elements each. 
The total number of spectral elements is $230,400$ and we use an order $7$ interpolating polynomial along each element so that the total resolution is $230,400\times8^3 \approx 1.2\times10^{8}$ points. 
Radial resolution was chosen such that at least one element spans the Ekman boundary layer.


\subsection{Topography}


\begin{figure*}
	\begin{center}
         \subfloat[][]{\includegraphics[width=0.3\textwidth]{Oliver_2026/figures/Y0201_topo}}
	\subfloat[][]{\includegraphics[width=0.3\textwidth]{Oliver_2026/figures/Y0804_topo}} 
	\subfloat[][]{\includegraphics[width=0.3\textwidth]{Oliver_2026/figures/Y3216_topo}}\\
	\subfloat[][]{\includegraphics[width=0.3\textwidth]{Oliver_2026/figures/Y3217_topo}} 
	\subfloat[][]{\includegraphics[width=0.3\textwidth]{Oliver_2026/figures/tomo_topo}}
	\subfloat[][]{\includegraphics[width=0.3\textwidth]{Oliver_2026/figures/hump_vis_white}}\\
	\end{center}
	\caption[Topographic models]{Visualisations of the topography used in the simulations: (a) $\hlm{2}{1}$;  (b) $\hlm{8}{4}$; (c) $\hlm{32}{16}$; (d)$\hlm{32}{17}$; (e) $h_t$. (f) The localized topography from \citep{tO25}.
	}

	\label{f:topo}
\end{figure*}


One of the primary goals in the present work is to better understand the impact of topography morphology on outer core convection.  
For most of the cases we therefore choose topographic shapes characterized by single spherical harmonics so that topographic length scales are easily defined. Thus, the topographic shape function corresponding to a spherical harmonic of degree $\ell$ and order $m$ is denoted as $\hlm{\ell }{m}$ and defined by
\begin{equation}
	\hlm{\ell}{m}\lb \theta,\phi\rb = \frac{1}{Q_{\ell}^{m}}P_{\ell}^{m} \lb \cos{\theta}\rb\; \cos\lb m \phi\rb,
	\label{e:hlm}
\end{equation}
where $P_{\ell}^{m}\lb \cos\theta\rb $ is the associated Legendre polynomial for degree $\ell$ and order $m,$ and $Q_{\ell}^{m}$ is the maximum value of $P_{\ell}^{m}\lb \cos\theta\rb .$
We study topographic amplitudes ranging from $0.01<\epsilon<0.2.$ All studied values of $\epsilon$ are larger than the estimated core value $\epsilon \approx 10^{-3},$ because we desire topography with amplitudes greater than the viscous Ekman boundary layer thickness $\delta_{E}=O\lb Ek^{1/2}\rb.$ In the core, a conservative estimate is $\epsilon /\delta_{E}=10^{4}$, and it is necessary to maintain this hierarchy of length scales. $\delta_{E}= O\lb 10^{-3}\rb $ in our simulations, so we study topographic amplitudes $\epsilon >\delta_{E}.$
We consider $\hlm{32}{16},\hlm{8}{4},$ and $\hlm{2}{1},$ across a sweep of $\epsilon$ and $\hlm{32}{17}$ for $\epsilon  = 0.1$.


We also study a topographic shape informed by seismological investigations of the deep mantle. 
%suggesting the presence of LLSVP's at the base of the mantle. 
For this shape, which we refer to as the ``tomographic model'' and denote the corresponding shape function as $h_{t}$, we use the cluster analysis maps developed in Lekic et al. (2012) \citep{vL12}, and consider only the areas where all of the included studies \citep{cH08,jR11,bK08,cM00,nS10} agree on the presence of an LLSVP.
We associate LLSVPs with intrusions into the core due to their high density \citep{jT15}, although this model is largely ad-hoc. 
Although CMB topography is likely broadband in structure, the voting map is essentially binary -- it shows where models agree and where they disagree concerning LLSVPs, but it does not inform us of the shape of the topography. Thus, a choice must be made with regards to representing this voting map. 
To smooth the edges which are otherwise sharp due to the discrete cluster analysis, we use an error function to build $h_{t}$. 
We fix the parameter determining the gradient of the topographic edges, so that $h_{t}$ depends on the value of the topographic amplitude according to
\begin{equation}
	h_{t}\lb \epsilon ,\theta,\phi\rb  = 0.5\lb1-\text{erf}\ls \delta \lb \theta,\phi\rb/\sigma \epsilon \rs \rb,
	\label{e:ht}
\end{equation}
where $\delta\lb \theta,\phi\rb $ gives the nearest distance to a cluster edge with the convention chosen such that $\delta\lb \theta,\phi\rb <0$ ($>0$) within (outside) the clusters. The slope parameter is $\sigma=28.2,$ which we choose because it corresponds to an average gradient of 0.02 per degree (that is a topography of amplitude $\epsilon=0.02$ decays to $10^{-2} \times\epsilon $ over $1^{\circ}$). 

Figures \ref{f:topo}(a)-(e) display the topographic shapes studied in the present work. For comparison, Figure \ref{f:topo}(f) displays the localized Gaussian bump investigated in Oliver et al. (2025).
In all cases $h\lb \theta,\phi\rb $ is normalized such that $\text{max}{\left|h\right|} = 1.0$. 
%such that $\epsilon$ can then be thought of as the max amplitude of the topography. 
We note that this is not the same as the orthonormal or  ``4$\pi$'' normalizations used in many texts, and that the surface area of the outer boundary is not constant across the different topographic shapes.
Since $\hlm{\ell}{m}$ is composed of a single harmonic, the surface area of the CMB is only weakly dependent on the topographic amplitude ($O\lb \epsilon^{2}\rb $). This is not the case for the tomographic model, where $h_{t}$ has a spherical mean. Table \ref{t:sa} gives the CMB surface areas for all geometries, showing that only small changes in the surface area are observed over our range of parameters.

During the discussion of topographic torques we present a scaling argument which requires a distribution rule for the size of the spherical harmonic coefficients of topography. 
Puica et al. (2023) \cite{mP23} and Dehant et al. (2025) \cite{vD25} used Kaula's rule \citep{wK66} to extend the CMB topography spectrum when calculating topographic torques, and we adopt the same approach here.
Kaula's rule suggests that the spherical harmonic coefficients for geophysical topography scale as $\ell^{-2}$. 
The tomography shape that we employ is artificial and so it is helpful to compare it with Kaula's rule.
Figure \ref{f:ht_spec} shows the magnitudes of the spectral coefficients for $\epsilon  = 0.20,$ $h_{t},$ where we observe that the spectrum roughly follows Kaula's rule. This rough agreement is interpreted as coincidental although it may be that the coarse cluster analysis performed in Lekic et al. (2012) \cite{vL12} is sufficient to reproduce the rule. 
Regardless, the plot in Figure \ref{f:ht_spec} does not reflect any physical result, it merely indicates that we expect our argument using Kaula's rule to apply to the tomography model.  It is straightforward to adjust the scaling law according to any topography spectrum.

\begin{figure}
	\begin{center}
		\subfloat[][]{\includegraphics[width=0.48\textwidth]{Oliver_2026/figures/ht_spec}}
	\end{center}
	\caption[Spectral decomposition of the tomographic shape function]{Spectral decomposition of the tomographic shape function, $h_{t}$ at $\epsilon  = 0.20$, averaged over $m$. The solid black line shows the scaling for Kaula's rule \citep{wK66}, which we use to estimate the torques. The dashed black line shows $\ell^{-1}$ for reference.}
	\label{f:ht_spec}
\end{figure}





\subsection{The geostrophic height function}


\begin{figure}
	\begin{center}
        \subfloat[][]{\includegraphics[width=0.19\textwidth]{Oliver_2026/figures/0x10Y0201_geo_contours}}
	\subfloat[][]{\includegraphics[width=0.19\textwidth]{Oliver_2026/figures/0x10Y0804_geo_contours}}
	\subfloat[][]{\includegraphics[width=0.19\textwidth]{Oliver_2026/figures/0x10Y3216_geo_contours}}
	\subfloat[][]{\includegraphics[width=0.19\textwidth]{Oliver_2026/figures/0x10Y3217_geo_contours}}
	\subfloat[][]{\includegraphics[width=0.19\textwidth]{Oliver_2026/figures/0x10tomo_geo_contours}} \\
	\end{center}
	\caption[Example gesotrophic contours]{Geostrophic contours and impact on time-averaged flow. Geostrophic contours for (left to right) $\hlm{2}{1}$, $\hlm{8}{4}$, $\hlm{32}{16}$, $\hlm{32}{17}$, and $h_{t}$ for $\epsilon  = 0.1$. }
	\label{f:t_geo_c}
\end{figure}


Using a cylindrical coordinate system ($s,\phi, z$), the geostrophic height function, $H_{geo}\lb s , \phi \rb$, measures the total axial height of the domain.  
%indicates the length of the chord which runs parallel to the rotation axis, intersects the point $\mathbf{r}$, and terminates at the boundaries of the solution domain. 
Curves of constant $H_{geo}$ are known as the geostrophic contours since motions that are in a purely geostrophic balance follow these contours \citep{hG68} (see Appendix \ref{a:oddveven} and Figure \ref{f:sketch} for a more complete discussion). 
As an example, $H_{geo}$ for a full sphere of radius $r_{o}$ is
\[H_{geo}^{s phere}\lb \mathbf{r}\rb =  H_{geo}^{s phere}\lb s\rb = 2\sqrt{r_{o}^{2} - s^{2}}.\]
In a spherical shell (inner radius $r_{i}$), $H_{geo}^{shell}$ has a discontinuity at the tangent cylinder $s = r_i$. 
%as the associated chord runs between the outer and inner boundary for  $s<r_{i}$ rather than between antipodal points on the outer surface. 
Nevertheless, the geostrophic contours are circles centered about the rotation axis.
The isosurfaces of $H_{geo}^{shell}$ are cylinders about the origin, and the term ``geostrophic cylinder'' commonly refers to these surfaces.
When topography is added to the outer spherical boundary the geostrophic contours are, in general, no longer circular, nor do they necessarily form closed curves. Sufficiently steep topography can introduce discontinuities in $H_{geo},$ which has significant consequences for the resulting flow. 
In this study we will refer to the isosurfaces of $H_{geo}$ as geostrophic cylinders as long as the associated contours form closed loops.
%We use the notation $H_{\ell,m}^{\delta},$ to refer to $H_{geo}$ for the $\epsilon =\delta,\hlm{\ell}{m}$ topography. \textcolor{red}{I don't understand this notation. Do you need the $\delta$ exponent?}

Figure \ref{f:t_geo_c} displays geostrophic contours for the topographic shapes investigated in the present work when $\epsilon  = 0.1$. 
Figures \ref{f:t_geo_c}(b), corresponding to $\hlm{32}{16}$, and \ref{f:t_geo_c}(c), $\hlm{32}{17}$, illustrate an important distinction between topographies where $\ell - m$ is even and odd.
In our notation $\hlm{\ell}{m} (\theta,\phi) = \lb -1\rb ^{\ell - m} \hlm{\ell}{m}\lb \pi-\theta,\phi\rb $, which indicates that when $\ell - m$ is even, the topography at antipodal points are identical. When $\ell - m$ is odd, the topography differs in sign.
An important consequence is that the perturbations to $H_{geo}$ are larger amplitude and have longer wavelength when $\ell-m $ is even. In particular, outside the tangent cylinder
\begin{equation}
	H_{geo}^{shell}- H_{\ell,m}^{\epsilon}\sim
	\begin{cases}
		 \epsilon  \cos m\phi\qquad &\ell - m \text{ even}\\
		 \epsilon^{2}  \cos 2m\phi\qquad &\ell - m \text{ odd.}
	\end{cases}
	%\begin{aligned}
	%\end{aligned}
	\label{e:h_geo}
\end{equation}
A derivation of this result is included in the Appendix \ref{a:oddveven}.

\subsection{Output/definitions}

To quantify flow speeds, we will report the Reynolds number $Re,$ which is calculated from the root mean square velocities according to
\be
Re = \sqrt{\left<\ub\cdot\ub\right>},
	\label{e:re}
	\ee
where the $\left< \cdot\right>$ indicates a full volume average. Unless otherwise noted, all reported values of $Re$ will also be time averaged. 
We also report the Reynolds number for particular components of the velocity field. 
We denote the included fields in the subscript so that (for example) $Re_{u,v}= \sqrt{\left<u^{2} + v^{2}\right>}$.
We discuss the geostrophic and ageostrophic Reynolds numbers ($Re_{\gchar}$ and $Re_{\achar}$ respectively) determined from the components of the velocity field tangent and perpendicular to the geostrophic contours. For these two quantities the average is performed over only the equatorial plane. All of the turbulent simulations are performed at the same value of $Ek$ and $\Rat,$ however the values of $Re$ vary. In the case with no topography (hereafter referred to as the control) $Re = 1280.$ Our most turbulent simulation ($\epsilon=  0.2, \;\hlm{8}{4}$) has a Reynolds number of $Re = 2480$.


The Nusselt number ($Nu$) compares the ratio of convective to conductive heat transport over a closed surface enclosing the inner core. 
The conductive heat profile differs between topographies, but the effect is negligible for small amplitudes. We therefore approximate the conductive heat profile using the temperature field that solves Laplace's equation in a sphere 
\be
\vartheta_{c}\lb r\rb  = \frac{r_{i}r_{o}}{ r} - r_{i}.
\label{e:tec}
\ee
%\[\nabla^{2}T_{c} = 0\qquad T_{c}\lb r_{i}\rb =1\qquad T_{c}\lb r_{o}\rb =0.\] 
 $Nu $ can be calculated over any surface, but for a statistically steady flow, the time averaged value is the same for any surface. 
We therefore present a Nusselt number calculated over the inner core boundary (ICB) surface according to
\[Nu = r_{i}\lb 1-r_{i}/r_{o}\rb\;\oint_{ICB}\frac{1}{r} \frac{\partial }{\partial r} \lb r\vartheta\rb da .\]


One dimensional (azimuthal) kinetic energy spectra are computed for various components of the flow in the equatorial plane. 
%For the velocity field, the calculation is sensitive to the chosen vector decomposition, so that it is necessary to indicate the scalar fields we break $\mathbf{u}$ into. 
In general, we define the radially averaged one dimensional energy spectra over fields $v_{i}$ as
\begin{equation}
	\mathcal{E}_{m}\lb v_{0},v_{1}...,v_{N}\rb = \lb r_c -r_{i}\rb ^{-1} 
	\int_{r_i}^{r_{c}} \sum_{i}^{N}\hat{v}_{i}^{m}\lb r\rb \hat{v}_{i}^{m^*}\lb r\rb  \partial r ,
\label{e:spectra_def}
\end{equation}
%We present an ageostrophic, two dimensional kinetic energy spectra of the flow in the equatorial plane. The ageostrophic energy spectral density is 
% \begin{equation}
%	 \mathcal{E}_{m} =  \left<\sum_{m}  \left| \mathcal{F}\lb\mathbf{u} - u_{g}\tvec\rb_{m}\right|^{2}\right>,
%	\label{e:em}
%\end{equation}
%where $\mathcal{F}(\chi)_{m}$ is the $m$ Fourier component of the field $\chi$ and $\left|\chi\right|$ indicates the 2-norm. 
where $\hat{\lb\cdot\rb}^{m} $ is the $m$ Fourier component of the field $\lb \cdot\rb $, and $r_c$ is the largest radial value within the equatorial plane for all $\phi$.    
The Fourier transform is performed over the azimuthal coordinate $\phi$. 
In particular, the kinetic energy spectra in a cylindrical coordinate system is $\mathcal{E}_{m} \lb u_{s},u_{\phi},w\rb.$  
We use the shorthand $\mathcal{E}_{m}=\mathcal{E}_{m}\lb  u_{s},u_{\phi},w\rb $ for notational simplicity. 

To quantify the flow morphology we report a weighted wavenumber $m_{eq}$, based on the kinetic energy spectra in the equatorial plane. It is defined as 
\begin{equation}
	m_{eq} = \frac{\sum_{m=1}m \mathcal{E}_{m}}{\sum_{m=1}\mathcal{E}_{m}}.
	\label{e:meq}
\end{equation}
The weighted wavenumber is similar to the mean spherical harmonic degree used in Christensen and Aubert (2006) \citep{uC06} (also \citep{jN24}), however it is defined for the order rather than the degree.
We found $m_{eq}$ was a useful quantity because the topography significantly alters the azimuthal structure of the flow.

%As a reminder, $\mathcal{E}_{m}$ and $m_{eq}$ are calculated from data in the equatorial plane only.


The topographic torque is defined as
\be
\mathbf{\Gamma} = \oint_{CMB} \mathbf{r}\times p \; d\mathbf{a},
\label{e:ttorque}
\ee
where $d\mathbf{a}$ is a surface area element on the CMB. We use the convention that the normal points away from the core. As in previous work \citep{tO25}, we will limit our analysis to the axial torque, which is responsible for LOD variations. For small topographic amplitude ($\ep \ll 1$) an approximate expression for the axial torque is 
\be
\Gamma_z \approx \oint \epsilon h\lb\theta,\phi\rb \partial_\phi p\; da.
\label{e:ax_torque}
\ee
All torques presented in this study are calculated with the exact expression (\ref{e:ttorque}), however we will appeal to equation (\ref{e:ax_torque}) for the purposes of analysis. 
The approximation in equation \eqref{e:ax_torque} is accurate up to $O\lb \epsilon^{2}\rb$ \citep{pR12}.

Motivated by equation \eqref{e:ax_torque}, we find it useful to investigate the spherical harmonic decomposition of the pressure field to better understand which components of the pressure contribute to the torque. For the $\hlm{\ell}{m}$ topographies, only a single harmonic contributes to the torque, whereas for the tomography model the entire spectrum of $p$ is convolved with $h_{t}.$ 
We report fluctuations in $\Gamma_{z}$ so rather than present the spectra of $p,$ we present the rms spectra of $p^{\prime}=p-\overline{p}$ on the outer surface. The overline $\overline{\cdot}$ denotes a time average.
We define the quantity
\[\phatLM = \sqrt{\overline {c_{\ell,m}^{2} + s_{\ell,m}^{2} } },\]
%\[p^{\prime m}_{\ell} = \sqrt{\overline {\left| p^{m}_{\ell}-\overline{p^{m}_{\ell}}\right|^{2}}},\]
where $c_{\ell,m}$ and $s_{\ell,m}$ satisfy
\[p(r_{cmb})-\overline{p}\lb r_{cmb}\rb = \sum_{\ell=0}\sum_{m = 0}^{\ell}P_{\ell}^{m}\lb \cos\theta\rb\; (c_{\ell,m}\cos(m\phi) + s_{\ell,m}\sin(m\phi)).\]
We refer $\phatLM$ as the fluctuating pressure spectrum.
For the $\hlm{\ell}{m}$ topographies, which are proportional to $\cos \lb m\phi\rb$, only the power in the corresponding $s_{\ell,m}$ mode contributes to the torque, so we find it useful to separate the sine and cosine modes. The data in Figure \ref{f:p_disp} reports 
$C_{\ell,m}=\sqrt{\overline{c_{\ell,m}^{2} } }$ and
$S_{\ell,m} =\sqrt{\overline{s_{\ell,m}^{2} } }$ separately.

We define the viscous torque as  
\[\mathbf{\Gamma}^{v} = \oint_{CMB} [(d\mathbf{a}\cdot\nabla) \mathbf{u} ]\times \mathbf{r},\]
although we only consider the axial component $\Gamma_{z}^{v}$. 
We report  the ratio of the standard deviations of the viscous and topographic axial torques
\[\gamma = \frac{\sqrt{\overline{(\Gamma_{z}^{v}-\overline{\Gamma_{z}^{v}})^{2}}}}{\sqrt{\overline{(\Gamma_{z}-\overline{\Gamma_{z}})^{2}}}}.\]
We found that $\gamma<1$ for all of the investigated parameters except the $\hlm{2}{1}$ topographies at $\epsilon \leq0.02.$ As in previous studies \citep[e.g.][]{bB07}, we conclude that viscous torques do not significantly contribute to $\Delta$LOD. Our results regarding $\gamma$ are provided in Appendix \ref{a:visc_t}.

